\documentclass[11pt, aspectratio=169]{beamer}

% ============================================================================
% PACKAGES ET CONFIGURATION
% ============================================================================

\usepackage[utf8]{inputenc}
\usepackage[french]{babel}
\usepackage{graphicx}
\usepackage{amsmath}
\usepackage{amssymb}
\usepackage{xcolor}
\usepackage{tikz}
\usepackage{listings}
\usepackage{fancybox}
\usepackage{hyperref}
\usepackage{multicol}
\usepackage{array}
\usepackage{booktabs}
\usepackage{lmodern}
\usepackage{fontawesome5}

% Configuration listings pour le code
\lstset{
    basicstyle=\ttfamily\small,
    breaklines=true,
    keywordstyle=\color{blue},
    commentstyle=\color{gray},
    stringstyle=\color{red},
    showstringspaces=false,
    breakatwhitespace=true,
    language=Java,
    framextopmargin=5pt,
    framebottom=5pt,
    frameround=tttt
}

% ============================================================================
% COULEURS PERSONNALISÉES - THÈME BLEU BLANC
% ============================================================================

\definecolor{DarkBlue}{RGB}{25, 80, 150}      % Bleu foncé principal
\definecolor{LightBlue}{RGB}{173, 216, 230}   % Bleu clair
\definecolor{MediumBlue}{RGB}{65, 130, 200}   % Bleu moyen
\definecolor{DarkGray}{RGB}{50, 50, 50}       % Gris foncé
\definecolor{LightGray}{RGB}{240, 240, 240}   % Gris clair
\definecolor{Accent}{RGB}{220, 100, 50}       % Orange accent
\definecolor{Success}{RGB}{76, 175, 80}       % Vert succès

% ============================================================================
% THÈME PERSONNALISÉ
% ============================================================================

\usetheme{Madrid}
\usecolortheme{default}

% Redéfinition des couleurs du thème
\setbeamercolor{primary}{bg=DarkBlue, fg=white}
\setbeamercolor{secondary}{bg=LightBlue, fg=DarkBlue}
\setbeamercolor{structure}{fg=DarkBlue}
\setbeamercolor{alerted text}{fg=Accent}
\setbeamercolor{example text}{fg=Success}

% Couleurs des titres et éléments
\setbeamercolor{title}{fg=white, bg=DarkBlue}
\setbeamercolor{subtitle}{fg=LightGray, bg=DarkBlue}
\setbeamercolor{frametitle}{fg=white, bg=MediumBlue}
\setbeamercolor{block title}{fg=white, bg=DarkBlue}
\setbeamercolor{block body}{fg=DarkGray, bg=LightGray}
\setbeamercolor{block title alerted}{fg=white, bg=Accent}
\setbeamercolor{block body alerted}{fg=DarkGray, bg=white}

% Style des boîtes
\setbeamertemplate{blocks}[rounded][shadow=true]

% Suppression du logo par défaut
\setbeamertemplate{sidebar right}{}
\setbeamertemplate{navigation symbols}{}

% Pied de page
\setbeamertemplate{footline}{%
    \leavevmode%
    \hbox{%
        \begin{beamercolorbox}[wd=.333333\paperwidth, ht=2.25ex, dp=1ex, center]{date in head/foot}%
            \usebeamerfont{date in head/foot}
            \insertframenumber{} / \inserttotalframenumber{}
        \end{beamercolorbox}%
        \begin{beamercolorbox}[wd=.333333\paperwidth, ht=2.25ex, dp=1ex, center]{date in head/foot}%
            \usebeamerfont{date in head/foot}
            Gestion Supermarché
        \end{beamercolorbox}%
        \begin{beamercolorbox}[wd=.333333\paperwidth, ht=2.25ex, dp=1ex, right]{date in head/foot}%
            \usebeamerfont{date in head/foot}
            \insertdate{}\hspace*{2em}
        \end{beamercolorbox}%
    }%
    \vskip0pt%
}

% ============================================================================
% INFORMATIONS DU DOCUMENT
% ============================================================================

\title[Gestion Supermarché]{Système de Gestion de Supermarché}
\subtitle{Projet Jakarta EE - Groupe 07}
\author{
    \large \textbf{NANA TCHOFFO STEVE JUNIOR} \\
    \normalsize Co-développeurs : [À compléter] \\
    \normalsize Université de Yaoundé 1
}
\institute{
    Filière : ICT4D \\
    Cours : ICT318 - Jakarta Enterprise Edition \\
    Année Académique : 2024-2025
}
\date{\today}

% ============================================================================
% DÉBUT DU DOCUMENT
% ============================================================================

\begin{document}

% ============================================================================
% PAGE DE GARDE
% ============================================================================

\begin{frame}[plain]
    \begin{tikzpicture}[remember picture, overlay]
        % Fond bleu gradient
        \fill[color=DarkBlue] (current page.southwest) rectangle (current page.northeast);
        
        % Bande blanche décorative
        \fill[color=white, opacity=0.1] (current page.northwest) -- ++(0, -2cm) -- ++(current page.width, 0) -- ++(0, 2cm) -- cycle;
        
        % Décoration géométrique
        \draw[color=white, opacity=0.15, line width=2pt] (2, 3) rectangle (5, 5);
        \draw[color=white, opacity=0.15, line width=2pt] (11, 1) circle (1.5cm);
        \draw[color=LightBlue, opacity=0.15, line width=2pt] (current page.east) circle (3cm);
    \end{tikzpicture}
    
    \begin{center}
        \vspace{2cm}
        
        % Logo placeholder (l'utilisateur peut insérer le logo)
        \framebox[1.2cm]{%
            \parbox[c][1.2cm]{1.2cm}{%
                \centering%
                \small LOGO UY1
            }%
        }
        
        \vspace{1.5cm}
        
        % Titre principal
        {\huge \textbf{\textcolor{white}{SYSTÈME DE GESTION DE SUPERMARCHÉ}}}
        
        \vspace{0.5cm}
        
        {\Large \textcolor{LightBlue}{Projet Jakarta Enterprise Edition}}
        
        \vspace{2cm}
        
        % Infos étudiant principal
        {\Large \textbf{\textcolor{white}{NANA TCHOFFO STEVE JUNIOR}}}
        
        \vspace{0.3cm}
        {\normalsize \textcolor{LightGray}{Étudiant principal - Développeur Backend}}
        
        \vspace{1.5cm}
        
        % Infos institution
        {\normalsize \textcolor{white}{
            \begin{tabular}{c}
                \textbf{Université de Yaoundé 1} \\
                Filière : ICT4D \\
                Cours : ICT318 \\
                Groupe : 07
            \end{tabular}
        }}
        
        \vspace{2cm}
        
        % Date
        {\small \textcolor{LightGray}{\today}}
    \end{center}
\end{frame}

% ============================================================================
% DIAPO 2 : TABLE DES MATIÈRES
% ============================================================================

\begin{frame}[plain]
    \frametitle{Table des Matières}
    
    \begin{multicols}{2}
        \begin{itemize}
            \item[\textcolor{DarkBlue}{\large \textbf{I.}}] Présentation du Projet
            \item[\textcolor{DarkBlue}{\large \textbf{II.}}] Objectives et Fonctionnalités
            \item[\textcolor{DarkBlue}{\large \textbf{III.}}] Architecture Générale
            \item[\textcolor{DarkBlue}{\large \textbf{IV.}}] Modèle de Données
            \item[\textcolor{DarkBlue}{\large \textbf{V.}}] Diagrammes UML
            \item[\textcolor{DarkBlue}{\large \textbf{VI.}}] Stack Technologique
            \item[\textcolor{DarkBlue}{\large \textbf{VII.}}] Configuration et Déploiement
            \item[\textcolor{DarkBlue}{\large \textbf{VIII.}}] Conclusion
        \end{itemize}
    \end{multicols}
\end{frame}

% ============================================================================
% SECTION I : PRÉSENTATION DU PROJET
% ============================================================================

\section{Présentation du Projet}

\begin{frame}
    \frametitle{I. Présentation du Projet}
    
    \begin{block}{Contexte}
        Développement d'une application web de gestion complète pour un supermarché,
        intégrant la gestion des produits, des ventes, des fournisseurs et des stocks.
    \end{block}
    
    \vspace{1cm}
    
    \begin{columns}[T]
        \column{0.45\textwidth}
        \begin{block}{Domaine Métier}
            \begin{itemize}
                \item[\faCheckCircle] Retail / Commerce
                \item[\faCheckCircle] Gestion d'Inventaire
                \item[\faCheckCircle] Point of Sale (POS)
                \item[\faCheckCircle] Supply Chain
            \end{itemize}
        \end{block}
        
        \column{0.45\textwidth}
        \begin{block}{Type d'Application}
            \begin{itemize}
                \item[\faServer] Web Application
                \item[\faCode] Backend : Jakarta EE
                \item[\faDisplay] Frontend : JSP + Bootstrap
                \item[\faDatabase] BD : MySQL 8
            \end{itemize}
        \end{block}
    \end{columns}
\end{frame}

% ============================================================================
% SECTION II : OBJECTIFS ET FONCTIONNALITÉS
% ============================================================================

\section{Objectifs et Fonctionnalités}

\begin{frame}
    \frametitle{II.1 Objectifs Globaux}
    
    \begin{enumerate}
        \item \textbf{Centraliser} la gestion des données du supermarché
        \item \textbf{Automatiser} les processus de vente et d'inventaire
        \item \textbf{Optimiser} la gestion des stocks et des alertes
        \item \textbf{Générer} des rapports et bons de commande
        \item \textbf{Sécuriser} l'accès aux données avec authentification
        \item \textbf{Intégrer} les services SMS et Email pour les notifications
    \end{enumerate}
    
    \vspace{1.5cm}
    
    \begin{alertblock}{Résultat Attendu}
        Une plateforme robuste, sécurisée et ergonomique pour les caissiers,
        gestionnaires de stock et administrateurs.
    \end{alertblock}
\end{frame}

\begin{frame}
    \frametitle{II.2 Fonctionnalités Principales}
    
    \begin{columns}[T]
        \column{0.48\textwidth}
        \begin{block}{Module Gestion Produits}
            \begin{itemize}
                \item[\faBox] CRUD complet (Créer, Lire, Modifier, Supprimer)
                \item[\faBarcode] Codes-barres
                \item[\faList] Catégorisation
                \item[\faPricetag] Gestion des prix
                \item[\faCalendar] Dates de péremption
                \item[\faWarehouse] Gestion des stocks
            \end{itemize}
        \end{block}
        
        \column{0.48\textwidth}
        \begin{block}{Module Ventes}
            \begin{itemize}
                \item[\faCashRegister] Point de Vente (Caisse)
                \item[\faTicketAlt] Génération de tickets
                \item[\faMoneyBillWave] Modes de paiement multiples
                \item[\faChartLine] Historique des ventes
                \item[\faFileInvoice] Facturation
            \end{itemize}
        \end{block}
    \end{columns}
\end{frame}

\begin{frame}
    \frametitle{II.3 Fonctionnalités Principales (Suite)}
    
    \begin{columns}[T]
        \column{0.48\textwidth}
        \begin{block}{Module Fournisseurs}
            \begin{itemize}
                \item[\faTruck] CRUD des fournisseurs
                \item[\faFileContract] Bons de commande
                \item[\faPDF] Export PDF
                \item[\faEnvelope] Envoi par Email
                \item[\faSms] Notification SMS
                \item[\faCalendarDays] Délai de livraison
            \end{itemize}
        \end{block}
        
        \column{0.48\textwidth}
        \begin{block}{Module Stock}
            \begin{itemize}
                \item[\faBell] Alertes de stock minimum
                \item[\faEye] Surveillance en temps réel
                \item[\faShieldAlt] Listener pour alertes automatiques
                \item[\faFileExcel] Export CSV
                \item[\faChartBar] Dashboard statistiques
            \end{itemize}
        \end{block}
    \end{columns}
\end{frame}

\begin{frame}
    \frametitle{II.4 Fonctionnalités Principales (Administration)}
    
    \begin{block}{Module Authentification \& Sécurité}
        \begin{itemize}
            \item[\faLock] Authentification par login/mot de passe
            \item[\faCryptocurrency] Hashage BCrypt des mots de passe
            \item[\faShield] Filtres de sécurité (AuthFilter)
            \item[\faUserTie] Gestion des rôles (ADMIN, CAISSIER)
            \item[\faUserCheck] Contrôle d'accès basé sur les rôles (RBAC)
        \end{itemize}
    \end{block}
    
    \vspace{0.8cm}
    
    \begin{block}{Module Administration}
        \begin{itemize}
            \item[\faUsers] Gestion des utilisateurs
            \item[\faChartPie] Dashboard administrateur
            \item[\faSync] Synchronisation des données
            \item[\faToolbox] Maintenance et logs
        \end{itemize}
    \end{block}
\end{frame}

% ============================================================================
% SECTION III : ARCHITECTURE GÉNÉRALE
% ============================================================================

\section{Architecture Générale}

\begin{frame}
    \frametitle{III.1 Architecture Générale du Système}
    
    \begin{center}
        \begin{tikzpicture}[scale=0.9]
            % Presentation Layer
            \draw[fill=LightBlue, draw=DarkBlue, line width=2pt] (0, 8) rectangle (6, 9.5);
            \node at (3, 8.75) {\textbf{\large Presentation Layer}};
            \node[below] at (1.5, 8.3) {\small JSP};
            \node[below] at (3, 8.3) {\small HTML/CSS};
            \node[below] at (4.5, 8.3) {\small Bootstrap};
            
            % Business Layer
            \draw[fill=white, draw=DarkBlue, line width=2pt] (0, 5.5) rectangle (6, 7.5);
            \node at (3, 7.2) {\textbf{\large Business Layer}};
            \node[below] at (1.5, 6.5) {\small Servlets};
            \node[below] at (3, 6.5) {\small Business Logic};
            \node[below] at (4.5, 6.5) {\small Controllers};
            
            % Data Layer
            \draw[fill=LightGray, draw=DarkBlue, line width=2pt] (0, 3) rectangle (6, 5);
            \node at (3, 4.3) {\textbf{\large Data Access Layer}};
            \node[below] at (1.5, 3.6) {\small DAOs};
            \node[below] at (3, 3.6) {\small JDBC};
            \node[below] at (4.5, 3.6) {\small SQL Queries};
            
            % Database
            \draw[fill=Accent, draw=DarkBlue, line width=2pt] (0, 0.5) rectangle (6, 2.5);
            \node at (3, 1.8) {\textbf{\large Database Layer}};
            \node[below] at (3, 1.1) {\textbf{MySQL 8}};
            
            % Arrows
            \draw[->, line width=2pt, DarkBlue] (3, 8) -- (3, 7.5);
            \draw[->, line width=2pt, DarkBlue] (3, 5.5) -- (3, 5);
            \draw[->, line width=2pt, DarkBlue] (3, 3) -- (3, 2.5);
        \end{tikzpicture}
    \end{center}
    
    \vspace{0.5cm}
    
    \begin{small}
        \textbf{Architecture en 4 couches :} Présentation → Métier → Données → Base de Données
    \end{small}
\end{frame}

\begin{frame}
    \frametitle{III.2 Patrons de Conception}
    
    \begin{columns}[T]
        \column{0.48\textwidth}
        \begin{block}{MVC Pattern}
            \begin{itemize}
                \item \textbf{Model :} Entités (Produit, Vente, etc.)
                \item \textbf{View :} JSP avec JSTL
                \item \textbf{Controller :} Servlets Jakarta
            \end{itemize}
        \end{block}
        
        \column{0.48\textwidth}
        \begin{block}{DAO Pattern}
            \begin{itemize}
                \item Isoler la logique d'accès aux données
                \item Faciliter les tests unitaires
                \item Classes DAO pour chaque entité
            \end{itemize}
        \end{block}
    \end{columns}
    
    \vspace{1cm}
    
    \begin{columns}[T]
        \column{0.48\textwidth}
        \begin{block}{Singleton Pattern}
            \begin{itemize}
                \item Configuration centralisée (AppConfig)
                \item Connexion à la BD unique
            \end{itemize}
        \end{block}
        
        \column{0.48\textwidth}
        \begin{block}{Observer Pattern}
            \begin{itemize}
                \item Listeners pour alertes de stock
                \item Événements automatiques
            \end{itemize}
        \end{block}
    \end{columns}
\end{frame}

\begin{frame}
    \frametitle{III.3 Structure des Packages}
    
    \begin{small}
        \begin{itemize}
            \item \textbf{com.supermarche.entity} \quad Modèles métier
            \begin{itemize}
                \item \textit{Produit, Fournisseur, Utilisateur, Vente, LigneVente}
            \end{itemize}
            
            \item \textbf{com.supermarche.dao} \quad Accès aux données
            \begin{itemize}
                \item \textit{ProduitDAO, FournisseurDAO, UtilisateurDAO, VenteDAO, LigneVenteDAO}
            \end{itemize}
            
            \item \textbf{com.supermarche.servlet} \quad Contrôleurs
            \begin{itemize}
                \item \textit{HomeServlet, auth, dashboard, produit, fournisseur, vente, stock, stats}
            \end{itemize}
            
            \item \textbf{com.supermarche.filter} \quad Sécurité
            \begin{itemize}
                \item \textit{AuthFilter pour l'authentification}
            \end{itemize}
            
            \item \textbf{com.supermarche.listener} \quad Événements
            \begin{itemize}
                \item \textit{StockAlertListener pour les alertes}
            \end{itemize}
            
            \item \textbf{com.supermarche.util} \quad Utilitaires
            \begin{itemize}
                \item \textit{AppConfig, BCryptUtil, PDF, CSV, SMS, Email}
            \end{itemize}
        \end{itemize}
    \end{small}
\end{frame}

% ============================================================================
% SECTION IV : MODÈLE DE DONNÉES
% ============================================================================

\section{Modèle de Données}

\begin{frame}
    \frametitle{IV.1 Diagramme Entité-Relation (MER)}
    
    \begin{center}
        \begin{tikzpicture}[scale=0.85, node distance=3cm]
            % Table Utilisateur
            \draw[fill=LightBlue, draw=DarkBlue, line width=1.5pt] (1, 7) rectangle (4, 9);
            \node[anchor=north] at (2.5, 8.85) {\textbf{UTILISATEUR}};
            \draw (1, 8.7) -- (4, 8.7);
            \node[anchor=west] at (1.1, 8.4) {id (PK)};
            \node[anchor=west] at (1.1, 8.0) {nom, prénom};
            \node[anchor=west] at (1.1, 7.6) {login (U)};
            \node[anchor=west] at (1.1, 7.2) {mot\_de\_passe};
            
            % Table Fournisseur
            \draw[fill=LightBlue, draw=DarkBlue, line width=1.5pt] (6, 7) rectangle (9, 9);
            \node[anchor=north] at (7.5, 8.85) {\textbf{FOURNISSEUR}};
            \draw (6, 8.7) -- (9, 8.7);
            \node[anchor=west] at (6.1, 8.4) {id (PK)};
            \node[anchor=west] at (6.1, 8.0) {raison\_sociale};
            \node[anchor=west] at (6.1, 7.6) {téléphone, email};
            \node[anchor=west] at (6.1, 7.2) {délai\_livraison};
            
            % Table Produit
            \draw[fill=LightBlue, draw=DarkBlue, line width=1.5pt] (3, 3.5) rectangle (7, 5.5);
            \node[anchor=north] at (5, 5.35) {\textbf{PRODUIT}};
            \draw (3, 5.2) -- (7, 5.2);
            \node[anchor=west] at (3.1, 4.9) {id (PK)};
            \node[anchor=west] at (3.1, 4.5) {code\_barre (U)};
            \node[anchor=west] at (3.1, 4.1) {designation, catégorie};
            \node[anchor=west] at (3.1, 3.7) {prix\_achat, prix\_vente};
            
            % Table Vente
            \draw[fill=LightBlue, draw=DarkBlue, line width=1.5pt] (0.5, 0.5) rectangle (3.5, 2.5);
            \node[anchor=north] at (2, 2.35) {\textbf{VENTE}};
            \draw (0.5, 2.2) -- (3.5, 2.2);
            \node[anchor=west] at (0.6, 1.9) {id (PK)};
            \node[anchor=west] at (0.6, 1.5) {caissier\_id (FK)};
            \node[anchor=west] at (0.6, 1.1) {date\_heure};
            \node[anchor=west] at (0.6, 0.7) {total\_ttc};
            
            % Table LigneVente
            \draw[fill=LightBlue, draw=DarkBlue, line width=1.5pt] (5, 0.5) rectangle (8, 2.5);
            \node[anchor=north] at (6.5, 2.35) {\textbf{LIGNE\_VENTE}};
            \draw (5, 2.2) -- (8, 2.2);
            \node[anchor=west] at (5.1, 1.9) {id (PK)};
            \node[anchor=west] at (5.1, 1.5) {vente\_id (FK)};
            \node[anchor=west] at (5.1, 1.1) {produit\_id (FK)};
            \node[anchor=west] at (5.1, 0.7) {quantité, sous-total};
            
            % Relations
            \draw[->] (2, 7) -- (2, 5.5); % Caissier -> Vente
            \draw[->] (6, 7) -- (6, 5.5); % Fournisseur -> Produit
            \draw[->] (2, 2.5) -- (6, 2.5); % Vente -> LigneVente
            \draw[->] (6, 2.5) -- (6, 3.5); % LigneVente -> Produit
        \end{tikzpicture}
    \end{center}
\end{frame}

\begin{frame}
    \frametitle{IV.2 Entité Produit}
    
    \begin{block}{Classe Produit}
        \begin{lstlisting}[language=Java]
public class Produit {
    private Long id;
    private String codeBarre;           // Unique
    private String designation;
    private String categorie;
    private BigDecimal prixAchat;
    private BigDecimal prixVente;
    private int stockActuel;            // Quantité en stock
    private int stockMinimum;           // Seuil d'alerte
    private LocalDate datePeremption;
    private Long fournisseurId;         // FK -> Fournisseur
    private Fournisseur fournisseur;    // Référence objet
}
        \end{lstlisting}
    \end{block}
    
    \vspace{0.5cm}
    
    \begin{small}
        \textbf{Caractéristiques :} Code-barre unique, Gestion des stocks avec seuil minimum,
        Dates de péremption, Lien avec fournisseur
    \end{small}
\end{frame}

\begin{frame}
    \frametitle{IV.3 Entité Vente}
    
    \begin{block}{Classe Vente}
        \begin{lstlisting}[language=Java]
public class Vente {
    public enum ModePaiement {
        ESPECES, CARTE, MOBILE_MONEY
    }
    
    private Long id;
    private Long caissierId;            // FK -> Utilisateur
    private LocalDateTime dateHeure;
    private BigDecimal totalTtc;
    private ModePaiement modePaiement;
    private String numeroTicket;        // Unique
    private Utilisateur caissier;
    private List<LigneVente> lignes;    // Produits vendus
    
    public void calculerTotal() {
        // Somme des sous-totaux
    }
}
        \end{lstlisting}
    \end{block}
\end{frame}

\begin{frame}
    \frametitle{IV.4 Entité Utilisateur \& Fournisseur}
    
    \begin{columns}[T]
        \column{0.48\textwidth}
        \begin{block}{Utilisateur}
            \begin{lstlisting}[language=Java, basicstyle=\ttfamily\tiny]
public class Utilisateur {
    public enum Role {
        ADMIN, CAISSIER
    }
    
    private Long id;
    private String nom;
    private String prenom;
    private String login;  // Unique
    private String motDePasse;  // BCrypt
    private Role role;
    private boolean actif;
    private LocalDateTime dateCreation;
}
            \end{lstlisting}
        \end{block}
        
        \column{0.48\textwidth}
        \begin{block}{Fournisseur}
            \begin{lstlisting}[language=Java, basicstyle=\ttfamily\tiny]
public class Fournisseur {
    private Long id;
    private String raisonSociale;  // Unique
    private String telephone;
    private String email;
    private int delaiLivraisonJours;
    private String conditionsPaiement;
}
            \end{lstlisting}
        \end{block}
    \end{columns}
\end{frame}

% ============================================================================
% SECTION V : DIAGRAMMES UML
% ============================================================================

\section{Diagrammes UML}

\begin{frame}
    \frametitle{V.1 Cas d'Usage (Use Cases) - Caissier}
    
    \begin{center}
        \begin{tikzpicture}[scale=1.0]
            % Acteur
            \draw (1, 2) circle (0.3);
            \draw (1, 2) -- (1, 1.2);
            \draw (0.7, 1.5) -- (1.3, 1.5);
            \draw (0.7, 1.2) -- (1, 0.7);
            \draw (1.3, 1.2) -- (1, 0.7);
            \node[below] at (1, 0.5) {\small Caissier};
            
            % Système
            \draw[DarkBlue, line width=2pt] (3, 0.5) rectangle (6, 4);
            \node at (4.5, 3.7) {\textbf{Système}};
            
            % Use cases
            \draw[rounded rectangle, rounded rectangle arc=90, draw=DarkBlue, fill=LightBlue, line width=1.5pt] 
                (3.5, 3) rectangle (5.5, 3.5);
            \node at (4.5, 3.25) {Effectuer une vente};
            
            \draw[rounded rectangle, rounded rectangle arc=90, draw=DarkBlue, fill=LightBlue, line width=1.5pt] 
                (3.5, 2.2) rectangle (5.5, 2.7);
            \node at (4.5, 2.45) {Consulter stock};
            
            \draw[rounded rectangle, rounded rectangle arc=90, draw=DarkBlue, fill=LightBlue, line width=1.5pt] 
                (3.5, 1.4) rectangle (5.5, 1.9);
            \node at (4.5, 1.65) {Imprimer ticket};
            
            % Relations
            \draw[->] (1.3, 2) -- (3.5, 3.2);
            \draw[->] (1.3, 2) -- (3.5, 2.45);
            \draw[->] (1.3, 2) -- (3.5, 1.65);
        \end{tikzpicture}
    \end{center}
    
    \vspace{0.8cm}
    
    \textbf{Cas d'usage principaux :}
    \begin{itemize}
        \item Effectuer une vente avec liste de produits
        \item Consulter stock disponible
        \item Générer et imprimer ticket
    \end{itemize}
\end{frame}

\begin{frame}
    \frametitle{V.2 Cas d'Usage (Use Cases) - Administrateur}
    
    \begin{center}
        \begin{tikzpicture}[scale=0.95]
            % Acteur Admin
            \draw (0.8, 2) circle (0.3);
            \draw (0.8, 2) -- (0.8, 1.2);
            \draw (0.5, 1.5) -- (1.1, 1.5);
            \draw (0.5, 1.2) -- (0.8, 0.7);
            \draw (1.1, 1.2) -- (0.8, 0.7);
            \node[below] at (0.8, 0.4) {\small Admin};
            
            % Système
            \draw[DarkBlue, line width=2pt] (2.5, 0.3) rectangle (5.5, 4);
            \node at (4, 3.8) {\textbf{Système}};
            
            % Use cases
            \draw[rounded rectangle, rounded rectangle arc=90, draw=DarkBlue, fill=LightBlue, line width=1.5pt] 
                (3, 3.3) rectangle (5, 3.8);
            \node at (4, 3.55) {Gérer produits};
            
            \draw[rounded rectangle, rounded rectangle arc=90, draw=DarkBlue, fill=LightBlue, line width=1.5pt] 
                (3, 2.5) rectangle (5, 3);
            \node at (4, 2.75) {Gérer fournisseurs};
            
            \draw[rounded rectangle, rounded rectangle arc=90, draw=DarkBlue, fill=LightBlue, line width=1.5pt] 
                (3, 1.7) rectangle (5, 2.2);
            \node at (4, 1.95) {Générer bon de commande};
            
            \draw[rounded rectangle, rounded rectangle arc=90, draw=DarkBlue, fill=LightBlue, line width=1.5pt] 
                (3, 0.9) rectangle (5, 1.4);
            \node at (4, 1.15) {Voir statistiques};
            
            % Relations
            \draw[->] (1.1, 2) -- (3, 3.55);
            \draw[->] (1.1, 2) -- (3, 2.75);
            \draw[->] (1.1, 2) -- (3, 1.95);
            \draw[->] (1.1, 2) -- (3, 1.15);
        \end{tikzpicture}
    \end{center}
    
    \vspace{0.8cm}
    
    \textbf{Cas d'usage Admin :}
    \begin{itemize}
        \item CRUD complet produits et fournisseurs
        \item Génération de bons de commande PDF
        \item Dashboard statistiques
    \end{itemize}
\end{frame}

\begin{frame}
    \frametitle{V.3 Diagramme de Flux - Processus de Vente}
    
    \begin{center}
        \begin{tikzpicture}[scale=0.85, node distance=2cm]
            % Start
            \draw[fill=Success, circle, draw=DarkBlue, line width=1.5pt] (3, 10) circle (0.4);
            \node[white] at (3, 10) {S};
            
            % Steps
            \node[rectangle, rounded corners, draw=DarkBlue, fill=LightBlue, line width=1.5pt, minimum width=3cm] 
                (a) at (3, 8.5) {Sélectionner produits};
            
            \node[rectangle, rounded corners, draw=DarkBlue, fill=LightBlue, line width=1.5pt, minimum width=3cm] 
                (b) at (3, 7) {Ajouter à panier};
            
            \node[rectangle, rounded corners, draw=DarkBlue, fill=LightBlue, line width=1.5pt, minimum width=3cm] 
                (c) at (3, 5.5) {Vérifier stock};
            
            \node[rectangle, rounded corners, draw=DarkBlue, fill=LightBlue, line width=1.5pt, minimum width=3cm] 
                (d) at (3, 4) {Calculer total};
            
            \node[rectangle, rounded corners, draw=DarkBlue, fill=LightBlue, line width=1.5pt, minimum width=3cm] 
                (e) at (3, 2.5) {Choisir mode paiement};
            
            \node[rectangle, rounded corners, draw=DarkBlue, fill=LightBlue, line width=1.5pt, minimum width=3cm] 
                (f) at (3, 1) {Valider \& Imprimer};
            
            % End
            \draw[fill=Accent, circle, draw=DarkBlue, line width=1.5pt] (3, -0.5) circle (0.4);
            \node[white] at (3, -0.5) {F};
            
            % Arrows
            \draw[->, line width=1.5pt, DarkBlue] (3, 9.6) -- (a);
            \draw[->, line width=1.5pt, DarkBlue] (a) -- (b);
            \draw[->, line width=1.5pt, DarkBlue] (b) -- (c);
            \draw[->, line width=1.5pt, DarkBlue] (c) -- (d);
            \draw[->, line width=1.5pt, DarkBlue] (d) -- (e);
            \draw[->, line width=1.5pt, DarkBlue] (e) -- (f);
            \draw[->, line width=1.5pt, DarkBlue] (f) -- (3, -0.1);
        \end{tikzpicture}
    \end{center}
    
    \vspace{0.5cm}
    \begin{small}
        \textbf{Flux complet d'une transaction de vente}
    \end{small}
\end{frame}

\begin{frame}
    \frametitle{V.4 Diagramme de Classe Simplifié}
    
    \begin{small}
        \begin{center}
            \begin{tikzpicture}[scale=0.8]
                % Classe Produit
                \draw[draw=DarkBlue, fill=LightBlue, line width=1.5pt] (0, 6) rectangle (3, 8);
                \node[anchor=north] at (1.5, 7.85) {\textbf{Produit}};
                \draw (0, 7.7) -- (3, 7.7);
                \node[anchor=west, font=\tiny] at (0.1, 7.4) {- codeBarre: String};
                \node[anchor=west, font=\tiny] at (0.1, 7.1) {- prix\_achat: BigDecimal};
                \node[anchor=west, font=\tiny] at (0.1, 6.8) {- stockActuel: int};
                \draw (0, 6.5) -- (3, 6.5);
                \node[anchor=west, font=\tiny] at (0.1, 6.2) {+ getPrix()};
                
                % Classe Vente
                \draw[draw=DarkBlue, fill=LightBlue, line width=1.5pt] (4, 6) rectangle (7, 8);
                \node[anchor=north] at (5.5, 7.85) {\textbf{Vente}};
                \draw (4, 7.7) -- (7, 7.7);
                \node[anchor=west, font=\tiny] at (4.1, 7.4) {- numeroTicket: String};
                \node[anchor=west, font=\tiny] at (4.1, 7.1) {- totalTtc: BigDecimal};
                \node[anchor=west, font=\tiny] at (4.1, 6.8) {- lignes: List};
                \draw (4, 6.5) -- (7, 6.5);
                \node[anchor=west, font=\tiny] at (4.1, 6.2) {+ calculerTotal()};
                
                % Classe LigneVente
                \draw[draw=DarkBlue, fill=LightBlue, line width=1.5pt] (2, 3.5) rectangle (5, 5.5);
                \node[anchor=north] at (3.5, 5.35) {\textbf{LigneVente}};
                \draw (2, 5.2) -- (5, 5.2);
                \node[anchor=west, font=\tiny] at (2.1, 4.9) {- quantite: int};
                \node[anchor=west, font=\tiny] at (2.1, 4.6) {- sousTotal: BigDecimal};
                \node[anchor=west, font=\tiny] at (2.1, 4.3) {- produit: Produit};
                \draw (2, 4) -- (5, 4);
                \node[anchor=west, font=\tiny] at (2.1, 3.7) {+ calculer()};
                
                % Relations
                \draw[->] (3.5, 5.5) -- (1.5, 8);
                \draw[->] (3.5, 5.5) -- (5.5, 8);
            \end{tikzpicture}
        \end{center}
    \end{small}
\end{frame}

% ============================================================================
% SECTION VI : STACK TECHNOLOGIQUE
% ============================================================================

\section{Stack Technologique}

\begin{frame}
    \frametitle{VI.1 Technologies Backend}
    
    \begin{columns}[T]
        \column{0.48\textwidth}
        \begin{block}{Plateforme}
            \begin{itemize}
                \item[\faJava] \textbf{Java 17}
                \item[\faCode] \textbf{Jakarta EE 10}
                \item[\faDatabase] \textbf{MySQL 8.3}
            \end{itemize}
        \end{block}
        
        \column{0.48\textwidth}
        \begin{block}{Framework Web}
            \begin{itemize}
                \item[\faServer] \textbf{Servlets Jakarta}
                \item[\faFileAlt] \textbf{JSP (JavaServer Pages)}
                \item[\faTag] \textbf{JSTL 3.0}
            \end{itemize}
        \end{block}
    \end{columns}
    
    \vspace{1cm}
    
    \begin{columns}[T]
        \column{0.48\textwidth}
        \begin{block}{Sécurité \& Crypto}
            \begin{itemize}
                \item[\faLock] \textbf{BCrypt} (Hashage mots de passe)
                \item[\faShield] \textbf{AuthFilter} (Authentification)
                \item[\faUserLock] RBAC (Contrôle d'accès)
            \end{itemize}
        \end{block}
        
        \column{0.48\textwidth}
        \begin{block}{Base de Données}
            \begin{itemize}
                \item[\faDatabase] \textbf{MySQL Connector/J 8.3}
                \item[\faCodeBranch] \textbf{JDBC} (Accès)
                \item[\faKey] Clés primaires et étrangères
            \end{itemize}
        \end{block}
    \end{columns}
\end{frame}

\begin{frame}
    \frametitle{VI.2 Technologies Frontend \& Utilitaires}
    
    \begin{columns}[T]
        \column{0.48\textwidth}
        \begin{block}{Présentation}
            \begin{itemize}
                \item[\faFileCode] \textbf{HTML5 / CSS3}
                \item[\faCheckSquare] \textbf{Bootstrap 5}
                \item[\faGithub] \textbf{Font Awesome} (Icônes)
            \end{itemize}
        \end{block}
        
        \column{0.48\textwidth}
        \begin{block}{Génération de Documents}
            \begin{itemize}
                \item[\faPDF] \textbf{iText 8.0.3} (PDF)
                \item[\faFileCsv] \textbf{OpenCSV 5.9} (Export)
            \end{itemize}
        \end{block}
    \end{columns}
    
    \vspace{1cm}
    
    \begin{columns}[T]
        \column{0.48\textwidth}
        \begin{block}{Communication}
            \begin{itemize}
                \item[\faSms] \textbf{Twilio 11.0.3} (SMS)
                \item[\faEnvelope] \textbf{Jakarta Mail} (Email)
            \end{itemize}
        \end{block}
        
        \column{0.48\textwidth}
        \begin{block}{Logging \& Build}
            \begin{itemize}
                \item[\faFileAlt] \textbf{SLF4J / Logback} (Logs)
                \item[\faToolbox] \textbf{Maven 3.x} (Build)
                \item[\faServer] \textbf{Jetty 12.0.7} (Runtime)
            \end{itemize}
        \end{block}
    \end{columns}
\end{frame}

\begin{frame}
    \frametitle{VI.3 Dépendances Maven}
    
    \begin{small}
        \begin{block}{Dépendances Principales}
            \begin{itemize}
                \item \texttt{jakarta.platform:jakarta.jakartaee-web-api:10.0.0}
                \item \texttt{jakarta.servlet.jsp.jstl:jakarta.servlet.jsp.jstl-api:3.0.0}
                \item \texttt{com.mysql:mysql-connector-j:8.3.0}
                \item \texttt{at.favre.lib:bcrypt:0.10.2}
                \item \texttt{com.itextpdf:kernel,layout:8.0.3}
                \item \texttt{com.opencsv:opencsv:5.9}
                \item \texttt{com.twilio.sdk:twilio:11.0.3}
                \item \texttt{org.eclipse.angus:jakarta.mail:2.0.3}
                \item \texttt{org.slf4j:slf4j-api:2.0.12}
                \item \texttt{ch.qos.logback:logback-classic:1.4.14}
            \end{itemize}
        \end{block}
    \end{small}
\end{frame}

% ============================================================================
% SECTION VII : CONFIGURATION ET DÉPLOIEMENT
% ============================================================================

\section{Configuration et Déploiement}

\begin{frame}
    \frametitle{VII.1 Configuration de la Base de Données}
    
    \begin{block}{Schéma SQL - Tables Principales}
        \begin{lstlisting}[language=SQL, basicstyle=\ttfamily\tiny]
CREATE TABLE utilisateur (
    id BIGINT AUTO_INCREMENT PRIMARY KEY,
    nom VARCHAR(100) NOT NULL,
    login VARCHAR(80) NOT NULL UNIQUE,
    mot_de_passe VARCHAR(255) NOT NULL,    -- BCrypt hashed
    role ENUM('ADMIN','CAISSIER'),
    actif BOOLEAN DEFAULT TRUE,
    date_creation DATETIME DEFAULT NOW()
);

CREATE TABLE fournisseur (
    id BIGINT AUTO_INCREMENT PRIMARY KEY,
    raison_sociale VARCHAR(150) UNIQUE,
    email VARCHAR(150),
    delai_livraison_jours INT DEFAULT 7
);

CREATE TABLE produit (
    id BIGINT AUTO_INCREMENT PRIMARY KEY,
    code_barre VARCHAR(30) UNIQUE,
    designation VARCHAR(150),
    categorie VARCHAR(80),
    prix_achat DECIMAL(10,2),
    prix_vente DECIMAL(10,2),
    stock_actuel INT DEFAULT 0,
    stock_minimum INT DEFAULT 5,
    date_peremption DATE,
    fournisseur_id BIGINT,
    FOREIGN KEY (fournisseur_id) REFERENCES fournisseur(id)
);
        \end{lstlisting}
    \end{block}
\end{frame}

\begin{frame}
    \frametitle{VII.2 Configuration de l'Application}
    
    \begin{block}{Variables d'Environnement (app.properties / env)}
        \begin{lstlisting}[language=bash, basicstyle=\ttfamily\small]
# Database Configuration
DB_URL=jdbc:mysql://localhost:3306/supermarche
DB_USERNAME=root
DB_PASSWORD=your_password

# Twilio SMS Configuration
SMS_ENABLED=true
TWILIO_ACCOUNT_SID=ACxxxxxxxxxxxxxxxxxxxxxxxxxxxxxxxx
TWILIO_AUTH_TOKEN=xxxxxxxxxxxxxxxxxxxxxxxxxxxxxxxx
TWILIO_FROM_NUMBER=+237657786440

# Email SMTP Configuration
MAIL_ENABLED=true
MAIL_SMTP_HOST=smtp.gmail.com
MAIL_SMTP_PORT=587
MAIL_SMTP_USERNAME=your_email@gmail.com
MAIL_SMTP_PASSWORD=your_app_password
MAIL_SMTP_FROM=your_email@gmail.com
        \end{lstlisting}
    \end{block}
\end{frame}

\begin{frame}
    \frametitle{VII.3 Déploiement et Exécution}
    
    \begin{block}{Étape 1 : Créer la Base de Données}
        \begin{lstlisting}[language=bash, basicstyle=\ttfamily\small]
# Lancer MySQL et créer la BD
mysql -u root -p < src/main/resources/schema.sql
        \end{lstlisting}
    \end{block}
    
    \vspace{0.5cm}
    
    \begin{block}{Étape 2 : Compiler avec Maven}
        \begin{lstlisting}[language=bash, basicstyle=\ttfamily\small]
# Compiler et packager
mvn clean package
        \end{lstlisting}
    \end{block}
    
    \vspace{0.5cm}
    
    \begin{block}{Étape 3 : Lancer l'Application}
        \begin{lstlisting}[language=bash, basicstyle=\ttfamily\small]
# Démarrer Jetty
mvn org.eclipse.jetty.ee10:jetty-ee10-maven-plugin:12.0.7:run

# Accéder à : http://localhost:8080/supermarche
        \end{lstlisting}
    \end{block}
\end{frame}

\begin{frame}
    \frametitle{VII.4 Comptes par Défaut \& Configuration}
    
    \begin{block}{Utilisateurs par Défaut}
        \begin{itemize}
            \item \textbf{Admin :} \texttt{login: admin} / \texttt{mot de passe: admin123} (Rôle: ADMIN)
            \item \textbf{Caissière :} \texttt{login: sophie} / \texttt{mot de passe: admin123} (Rôle: CAISSIER)
        \end{itemize}
    \end{block}
    
    \vspace{1cm}
    
    \begin{block}{Sécurité}
        Les mots de passe sont stockés en \textbf{BCrypt}, jamais en clair.
        
        \begin{lstlisting}[language=SQL, basicstyle=\ttfamily\tiny]
-- Mettre à jour les mots de passe (s'ils ne sont pas corrects)
UPDATE utilisateur
SET mot_de_passe = '$2a$12$kaiE7qTm97vdXSCqBVPYW.WOuLOB0ieoxwcfyYC.fCiy7SxuhxXWu'
WHERE login IN ('admin', 'sophie');
        \end{lstlisting}
    \end{block}
\end{frame}

% ============================================================================
% SECTION VIII : FONCTIONNALITÉS DÉTAILLÉES
% ============================================================================

\section{Fonctionnalités Détaillées}

\begin{frame}
    \frametitle{VIII.1 Module Gestion des Produits}
    
    \begin{block}{Servlets Impliquées}
        \begin{itemize}
            \item[\faList] \textbf{ListeProduitServlet} → Affiche liste paginée
            \item[\faWrench] \textbf{FormProduitServlet} → Affiche formulaire
            \item[\faPlus] \textbf{AjouterProduitServlet} → Crée nouveau produit
            \item[\faEdit] \textbf{ModifierProduitServlet} → Modifie produit existant
            \item[\faTrash] \textbf{SupprimerProduitServlet} → Supprime produit
        \end{itemize}
    \end{block}
    
    \vspace{0.8cm}
    
    \begin{block}{Fonctionnalités}
        \begin{itemize}
            \item CRUD complet (Create, Read, Update, Delete)
            \item Code-barre unique par produit
            \item Gestion du stock (actuel vs. minimum)
            \item Prix d'achat et de vente distincts
            \item Catégorisation des produits
            \item Dates de péremption
            \item Association fournisseur
        \end{itemize}
    \end{block}
\end{frame}

\begin{frame}
    \frametitle{VIII.2 Module Gestion des Fournisseurs}
    
    \begin{block}{Servlets Impliquées}
        \begin{itemize}
            \item[\faList] \textbf{ListeFournisseurServlet} → Liste des fournisseurs
            \item[\faWrench] \textbf{FormFournisseurServlet} → Formulaire
            \item[\faPlus] \textbf{AjouterFournisseurServlet} → Ajouter
            \item[\faEdit] \textbf{ModifierFournisseurServlet} → Modifier
            \item[\faFileContract] \textbf{BonCommandeFournisseurServlet} → Bon de commande
        \end{itemize}
    \end{block}
    
    \vspace{0.8cm}
    
    \begin{block}{Spécificités}
        \begin{itemize}
            \item CRUD des fournisseurs
            \item Génération de bons de commande en PDF
            \item Envoi automatique par Email
            \item Notification SMS optionnelle
            \item Délai de livraison
            \item Conditions de paiement
        \end{itemize}
    \end{block}
\end{frame}

\begin{frame}
    \frametitle{VIII.3 Module Caisse \& Ventes}
    
    \begin{block}{Processus}
        \begin{itemize}
            \item[\faCashRegister] Sélection des produits par code-barre
            \item[\faShoppingCart] Panier dynamique
            \item[\faCalculator] Calcul automatique des totaux
            \item[\faMoneyBillWave] Choix du mode de paiement (Espèces, Carte, Mobile)
            \item[\faTicketAlt] Génération de ticket
        \end{itemize}
    \end{block}
    
    \vspace{0.8cm}
    
    \begin{block}{Servlets}
        \begin{itemize}
            \item \textbf{CaisseServlet} → Interface caisse
            \item \textbf{ValiderVenteServlet} → Validation transaction
            \item \textbf{TicketVenteServlet} → Impression ticket
        \end{itemize}
    \end{block}
\end{frame}

\begin{frame}
    \frametitle{VIII.4 Module Stock \& Alertes}
    
    \begin{block}{Fonctionnalités}
        \begin{itemize}
            \item[\faBell] Alertes automatiques quand stock \textless{} minimum
            \item[\faEye] Monitoring en temps réel
            \item[\faWarning] Listener pour événements de stock
            \item[\faShoppingCart] Décrément automatique lors de ventes
            \item[\faTruck] Suggestion de commande fournisseur
        \end{itemize}
    \end{block}
    
    \vspace{0.8cm}
    
    \begin{block}{Implémentation}
        \textbf{StockAlertListener} utilise le pattern Observer pour :
        \begin{itemize}
            \item Notifier les administrateurs
            \item Déclencher les bons de commande automatiquement
            \item Logger les événements critiques
        \end{itemize}
    \end{block}
\end{frame}

\begin{frame}
    \frametitle{VIII.5 Module Authentification \& Sécurité}
    
    \begin{block}{Mécanisme}
        \begin{enumerate}
            \item Utilisateur saisit login + mot de passe
            \item \textbf{BCrypt} vérifie le mot de passe
            \item Création de session Jakarta
            \item \textbf{AuthFilter} contrôle chaque requête
            \item \textbf{RBAC} : Admin vs Caissier
        \end{enumerate}
    \end{block}
    
    \vspace{0.8cm}
    
    \begin{block}{Sécurité}
        \begin{itemize}
            \item[\faLock] Mots de passe hashés avec BCrypt
            \item[\faShieldAlt] Pas de session = redirection login
            \item[\faUserLock] Rôles contrôlent l'accès
            \item[\faEyeSlash] Données sensibles jamais loggées
            \item[\faLockOpen] Session expire après inactivité
        \end{itemize}
    \end{block}
\end{frame}

\begin{frame}
    \frametitle{VIII.6 Utilitaires \& Services}
    
    \begin{columns}[T]
        \column{0.48\textwidth}
        \begin{block}{AppConfig}
            \begin{itemize}
                \item Singleton de configuration
                \item Variables d'environnement
                \item Propriétés applicatives
                \item Connexion BD centralisée
            \end{itemize}
        \end{block}
        
        \column{0.48\textwidth}
        \begin{block}{Services}
            \begin{itemize}
                \item \textbf{PDFGenerator} : Bons, factures
                \item \textbf{CSVExporter} : Export données
                \item \textbf{SMSService} : Twilio
                \item \textbf{EmailService} : Jakarta Mail
            \end{itemize}
        \end{block}
    \end{columns}
\end{frame}

% ============================================================================
% SECTION IX : MODÈLE DE DONNÉES GRAPHIQUE
% ============================================================================

\section{Modèle de Données}

\begin{frame}
    \frametitle{IX. Schéma Complet de la Base de Données}
    
    \begin{center}
        \small
        \begin{tabular}{|c|c|c|}
            \hline
            \rowcolor{LightBlue}
            \textbf{Table} & \textbf{Colonnes Principales} & \textbf{Contraintes} \\
            \hline
            \hline
            \textbf{utilisateur} & id, nom, login, mot\_de\_passe, role, actif & PK: id, U: login \\
            \hline
            \textbf{fournisseur} & id, raison\_sociale, email, téléphone, délai & PK: id, U: rais\_soc \\
            \hline
            \textbf{produit} & id, code\_barre, designation, prix\_achat, stock & PK: id, U: code\_barre \\
            & categorie, date\_peremption, fournisseur\_id & FK: fournisseur\_id \\
            \hline
            \textbf{vente} & id, caissier\_id, date\_heure, total\_ttc & PK: id, FK: caissier\_id \\
            & mode\_paiement, numero\_ticket & U: numero\_ticket \\
            \hline
            \textbf{ligne\_vente} & id, vente\_id, produit\_id, quantite & PK: id \\
            & sous\_total & FK: vente\_id, produit\_id \\
            \hline
        \end{tabular}
    \end{center}
    
    \vspace{0.8cm}
    \begin{small}
        \textbf{Légende :} PK = Clé Primaire, FK = Clé Étrangère, U = Unique
    \end{small}
\end{frame}

% ============================================================================
% SECTION X : AVANTAGES \& DÉFIS
% ============================================================================

\section{Avantages et Défis}

\begin{frame}
    \frametitle{X.1 Avantages du Projet}
    
    \begin{columns}[T]
        \column{0.48\textwidth}
        \begin{block}{Architecture}
            \begin{itemize}
                \item[\faCheckCircle] Séparation des responsabilités
                \item[\faCheckCircle] Code modulaire et maintenable
                \item[\faCheckCircle] Patterns de conception reconnus
                \item[\faCheckCircle] Facilité d'extension
            \end{itemize}
        \end{block}
        
        \column{0.48\textwidth}
        \begin{block}{Fonctionnalités}
            \begin{itemize}
                \item[\faCheckCircle] Système complet de point de vente
                \item[\faCheckCircle] Alertes intelligentes de stock
                \item[\faCheckCircle] Export PDF/CSV
                \item[\faCheckCircle] Notifications SMS/Email
            \end{itemize}
        \end{block}
    \end{columns}
    
    \vspace{1cm}
    
    \begin{block}{Sécurité \& Performance}
        \begin{itemize}
            \item[\faCheckCircle] Authentification robuste (BCrypt)
            \item[\faCheckCircle] RBAC pour contrôle d'accès
            \item[\faCheckCircle] Requêtes SQL optimisées
            \item[\faCheckCircle] Logging structuré
        \end{itemize}
    \end{block}
\end{frame}

\begin{frame}
    \frametitle{X.2 Défis Rencontrés \& Solutions}
    
    \begin{block}{Défi 1 : Gestion Complexe des Transactions}
        \textbf{Solution :} Utilisation de DAOs avec transactions JDBC appropriées
    \end{block}
    
    \vspace{0.5cm}
    
    \begin{block}{Défi 2 : Synchronisation Stock}
        \textbf{Solution :} Pattern Observer avec StockAlertListener
    \end{block}
    
    \vspace{0.5cm}
    
    \begin{block}{Défi 3 : Génération PDF/CSV}
        \textbf{Solution :} Bibliothèques iText et OpenCSV
    \end{block}
    
    \vspace{0.5cm}
    
    \begin{block}{Défi 4 : Configuration Environnement}
        \textbf{Solution :} Variables d'environnement + fichiers propriétés
    \end{block}
    
    \vspace{0.5cm}
    
    \begin{block}{Défi 5 : Performance Requêtes}
        \textbf{Solution :} Index sur clés étrangères, paginationquête optimisées
    \end{block}
\end{frame}

% ============================================================================
% SECTION XI : ÉVOLUTIONS FUTURES
% ============================================================================

\section{Évolutions Futures}

\begin{frame}
    \frametitle{XI. Évolutions Futures Possibles}
    
    \begin{block}{Court Terme (Prochaines versions)}
        \begin{itemize}
            \item[\faArrowRight] Rapports détaillés par période
            \item[\faArrowRight] Dashboard analytique avancé
            \item[\faArrowRight] Prévisions de stock (Machine Learning)
            \item[\faArrowRight] Intégration paiements en ligne
        \end{itemize}
    \end{block}
    
    \vspace{0.8cm}
    
    \begin{block}{Moyen Terme}
        \begin{itemize}
            \item[\faArrowRight] Application mobile (Android/iOS)
            \item[\faArrowRight] Synchronisation multi-magasins
            \item[\faArrowRight] API REST pour intégrations tiers
            \item[\faArrowRight] Audit trail complet
        \end{itemize}
    \end{block}
    
    \vspace{0.8cm}
    
    \begin{block}{Long Terme}
        \begin{itemize}
            \item[\faArrowRight] Migration vers microservices
            \item[\faArrowRight] Blockchain pour traçabilité
            \item[\faArrowRight] Système d'IA pour recommandations
        \end{itemize}
    \end{block}
\end{frame}

% ============================================================================
% CONCLUSION
% ============================================================================

\section{Conclusion}

\begin{frame}
    \frametitle{Conclusion}
    
    \begin{block}{Résumé}
        Notre application \textbf{Gestion Supermarché} démontre :
        \begin{itemize}
            \item[\faCheckCircle] Maîtrise de \textbf{Jakarta EE}
            \item[\faCheckCircle] Compréhension des \textbf{patterns d'architecture}
            \item[\faCheckCircle] Implémentation d'un système \textbf{sécurisé} et \textbf{scalable}
            \item[\faCheckCircle] Intégration de \textbf{services externes} (SMS, Email, PDF)
            \item[\faCheckCircle] Bonnes pratiques de \textbf{développement professionnel}
        \end{itemize}
    \end{block}
    
    \vspace{1cm}
    
    \begin{alertblock}{Points Clés}
        Ce projet illustre comment construire une application web robuste
        pour un domaine métier réel (Retail), avec une architecture maintenable
        et des fonctionnalités complètes.
    \end{alertblock}
\end{frame}

\begin{frame}[plain]
    \begin{center}
        \vspace{2cm}
        {\Huge \textbf{\textcolor{DarkBlue}{Merci !}}}
        
        \vspace{1.5cm}
        
        {\Large Questions ?}
        
        \vspace{2cm}
        
        {\small
            \textbf{Contact :} NANA TCHOFFO STEVE JUNIOR \\
            \textbf{Université :} Yaoundé 1 - ICT4D \\
            \textbf{Cours :} ICT318 - Jakarta EE
        }
        
        \vspace{2cm}
        
        {\footnotesize \textcolor{gray}{
            Groupe 07 - Année Académique 2024-2025 \\
            Projet : Système de Gestion de Supermarché
        }}
    \end{center}
\end{frame}

\end{document}
